% =====================================================================
% Reporte_seleccion_operadores_y_su_efecto.tex
% ---------------------------------------------------------------------
% Documento monolitico que documenta y compara las estrategias de
% seleccion de operadores desarrolladas en el proyecto MetaCARP.
% Compilable con pdflatex (>= 2 pases para resolver referencias).
% =====================================================================
\documentclass[11pt,a4paper]{article}

% ----- Codificacion e idioma -----
\usepackage[utf8]{inputenc}
\usepackage[T1]{fontenc}
\usepackage[spanish,es-nodecimaldot]{babel}

% ----- Matematicas -----
\usepackage{amsmath, amssymb, amsthm}
\usepackage{mathtools}

% ----- Tablas profesionales -----
\usepackage{booktabs}
\usepackage{array}
\usepackage{tabularx}

% ----- Pseudocodigo -----
\usepackage{algorithm}
\usepackage{algpseudocode}
\algrenewcommand\algorithmicrequire{\textbf{Entradas:}}
\algrenewcommand\algorithmicensure{\textbf{Devuelve:}}
\algrenewcommand\algorithmicif{\textbf{si}}
\algrenewcommand\algorithmicthen{\textbf{entonces}}
\algrenewcommand\algorithmicelse{\textbf{si no}}
\algrenewcommand\algorithmicend{\textbf{fin}}
\algrenewcommand\algorithmicwhile{\textbf{mientras}}
\algrenewcommand\algorithmicfor{\textbf{para}}
\algrenewcommand\algorithmicdo{\textbf{hacer}}
\algrenewcommand\algorithmicreturn{\textbf{devolver}}
\algrenewcommand\algorithmicrepeat{\textbf{repetir}}
\algrenewcommand\algorithmicuntil{\textbf{hasta}}

% ----- Hipervinculos y color -----
\usepackage{xcolor}
\usepackage[colorlinks=true, linkcolor=blue!60!black,
            citecolor=blue!60!black, urlcolor=blue!60!black]{hyperref}

% ----- Margenes y espacios -----
\usepackage[a4paper,margin=2.4cm]{geometry}
\setlength{\parskip}{4pt}

% ----- Cabeceras -----
\usepackage{fancyhdr}
\pagestyle{fancy}
\fancyhf{}
\fancyhead[L]{MetaCARP -- Seleccion de operadores}
\fancyhead[R]{\thepage}
\renewcommand{\headrulewidth}{0.4pt}

% ----- Comandos abreviados -----
\newcommand{\carp}{CARP}
\newcommand{\mh}{MH}
\newcommand{\meta}[1]{\textsc{#1}}
\newcommand{\op}[1]{\texttt{#1}}
\newcommand{\code}[1]{\texttt{#1}}
\DeclareMathOperator*{\argmin}{arg\,min}
\DeclareMathOperator*{\argmax}{arg\,max}

% =====================================================================
\title{\textbf{Seleccion de operadores y su efecto sobre las
metaheuristicas del proyecto MetaCARP}\\[4pt]
\large Reporte minucioso de experimentos
y comparativa cuantitativa}
\author{Alhely Gonzalez Luna\\\small \texttt{alelygl@gmail.com}}
\date{Mayo 2026}
% =====================================================================

\begin{document}
\maketitle

\begin{abstract}
\noindent Este reporte recoge, justifica y compara todas las estrategias
de \emph{seleccion de operadores de vecindario} desarrolladas en el
proyecto \textsc{MetaCARP} para resolver el \emph{Capacitated Arc
Routing Problem} (\carp). Se documentan en detalle (i) los nueve
operadores implementados, (ii) las cinco metaheuristicas que los
consumen, (iii) las cuatro estrategias de seleccion ensayadas
(uniforme, sesgo binario fluido, sesgo binario estricto y
\emph{Adaptive Operator Selection} con \emph{Probability Matching}) y
(iv) la variante reciente de \emph{Path Relinking} truncado. Se
discuten resultados experimentales sobre la suite estandar y se
extraen lecciones cualitativas para futuras extensiones.
\end{abstract}

\tableofcontents
\newpage

% ===================================================================
\section{Introduccion}
% ===================================================================

\subsection{Definicion formal del \carp}

El \emph{Capacitated Arc Routing Problem} (\carp) toma como entrada un
grafo no dirigido y conexo
\begin{equation}
G = (V, E),\qquad
E = E_R \cup E_{NR},\qquad E_R \cap E_{NR} = \emptyset,
\end{equation}
donde $V$ es el conjunto de vertices, $E_R$ el subconjunto de
\emph{aristas requeridas} (tareas) y $E_{NR}$ el de \emph{aristas no
requeridas} (transitos puros). Para cada arista $e \in E$ existen una
funcion de costo $c(e) \ge 0$ y una funcion de demanda
$d(e) \ge 0$ con $d(e) > 0$ si y solo si $e \in E_R$. Existe ademas un
nodo deposito $v_0 \in V$ y una flota de a lo sumo $K$ vehiculos
identicos con capacidad $Q$.

Una \emph{solucion} es una particion ordenada $S=(R_1,\dots,R_m)$ con
$m\le K$, donde cada $R_k$ es una secuencia de aristas requeridas
servidas por el vehiculo $k$, comenzando y terminando en $v_0$. El
problema consiste en encontrar
\begin{equation}\label{eq:carp}
\min_{S} \quad
\sum_{k=1}^{m} \bigl[\, \text{deadhead}(R_k) + \sum_{e\in R_k} c(e)\,\bigr]
\quad\text{sujeto a}\quad
\sum_{e\in R_k} d(e) \le Q,\quad
\bigcup_{k} R_k = E_R,
\end{equation}
con cada tarea servida exactamente una vez por un unico vehiculo. El
termino $\text{deadhead}(R_k)$ acumula los costos de los caminos mas
cortos en $G$ que enlazan dos servicios consecutivos.

\subsection{Por que la seleccion de operadores es critica}

Las metaheuristicas (en adelante, \mh) usadas en \textsc{MetaCARP}
son \emph{trayectoria} (Recocido Simulado, Tabu Search Simple,
Reactive Tabu Search) y \emph{poblacion} (\emph{Artificial Bee Colony},
\emph{Cuckoo Search}). Todas comparten la misma mecanica nuclear:
\emph{generar un vecino} de la solucion actual aplicando un operador
de vecindario, evaluarlo, y decidir si lo aceptan. Los operadores de
vecindario constituyen, por tanto, el \emph{unico} medio del que
dispone la \mh{} para moverse en el espacio de soluciones.

Existen dos categorias claramente distintas:
\begin{itemize}
\item \textbf{Operadores intra-ruta}: solo reordenan las tareas
\emph{dentro} de una ruta. No alteran la asignacion tarea $\to$ ruta
y, por consiguiente, \emph{no pueden reducir una violacion de
capacidad}.
\item \textbf{Operadores inter-ruta}: trasladan tareas \emph{entre}
rutas. Modifican la distribucion de demanda por vehiculo y son los
unicos capaces de reparar violaciones de capacidad.
\end{itemize}

Como la funcion objetivo agregada combina \emph{costo} y
\emph{violacion} (ver Seccion \ref{sec:representacion}), una eleccion
ciega de operadores acaba con la \mh{} eligiendo movimientos intra
cuando lo que necesita son inter, o viceversa. Por ello la
\emph{seleccion de operadores} es un eje de diseno tan importante como
los propios parametros termodinamicos, tabu o poblacionales.

\subsection{Roadmap del documento}

La Seccion \ref{sec:representacion} fija la notacion, el formato de
solucion y el objetivo penalizado. La Seccion \ref{sec:operadores}
documenta cada uno de los nueve operadores en detalle algoritmico. La
Seccion \ref{sec:metaheuristicas} hace lo propio con las cinco \mh{}.
La Seccion \ref{sec:seleccion} formaliza las cuatro estrategias de
seleccion estudiadas. La Seccion \ref{sec:experimentos} recoge los
experimentos en orden cronologico. La Seccion \ref{sec:comparativa}
presenta la tabla maestra de \emph{gap} al \emph{BKS}. La
Seccion \ref{sec:insights} discute las conclusiones cualitativas y la
Seccion \ref{sec:conclusiones} cierra con recomendaciones.

% ===================================================================
\section{Representacion de soluciones}\label{sec:representacion}
% ===================================================================

\subsection{Formato \texttt{list[list[str]]} (etiquetas)}

Cada solucion se representa internamente como una lista de rutas y
cada ruta como una lista de etiquetas string. Las tareas requeridas
reciben etiquetas \op{TR<n>} y las no requeridas \op{TNR<n>}. El
deposito se marca con la etiqueta \op{D} al inicio y al final de cada
ruta. Por ejemplo:
\begin{verbatim}
solucion = [["D", "TR1", "TR5", "TR3", "D"],
            ["D", "TR2", "TR4", "D"]]
\end{verbatim}

Este formato es la representacion \emph{canonica} para inspeccion
humana, reporte CSV y entrada/salida de funciones publicas.

\subsection{Formato \texttt{list[list[int]]} (IDs)}

Para evaluacion vectorizada y operacion sobre grandes lotes de
vecinos, \textsc{MetaCARP} provee una codificacion entera donde cada
tarea recibe un ID en $\{0,\dots,n-1\}$ y el deposito ID $-1$. Las
rutas no llevan deposito intercalado (los extremos son implicitos).
Por ejemplo el vector anterior se codifica como:
\begin{verbatim}
solucion_ids = [[0, 4, 2], [1, 3]]
\end{verbatim}

El \emph{Artificial Bee Colony} (ABC) usa este formato \emph{a lo
largo de toda la corrida} porque su bucle interno evalua decenas de
vecinos por iteracion y la conversion etiqueta $\to$ id en cada
evaluacion seria un cuello de botella; en contraste, las \mh{} de
trayectoria suelen operar en etiquetas y solo codifican a IDs para el
lote tabu.

\subsection{Marcador de deposito y normalizacion}

El simbolo \code{md\_op} (typicamente \op{"D"}) se obtiene del
contexto o se pasa explicitamente. Los operadores de vecindario
\emph{normalizan} la entrada eliminando los depots de los extremos
antes de aplicar la transformacion, y los reinsertan al devolver. Esto
desacopla la mecanica del operador de la convencion de marcado y
permite implementaciones compactas.

\subsection{Funcion objetivo penalizada}\label{sec:objetivo}

Sea $\text{costo}(S)$ la suma definida en \eqref{eq:carp} y
$\text{viol}(S) = \sum_k \max(0,\sum_{e\in R_k} d(e) - Q)$ el exceso
total de demanda. \textsc{MetaCARP} usa el \emph{objetivo penalizado}
\begin{equation}\label{eq:objetivo}
f_\lambda(S) \;=\; \text{costo}(S) \;+\; \lambda \cdot \text{viol}(S),
\qquad \lambda \ge 0.
\end{equation}

Esta funcion combina objetivos en un escalar y, con $\lambda$
suficientemente grande, garantiza que el optimo de
$f_\lambda$ sea siempre factible. Las cinco \mh{} comparan vecinos
sobre $f_\lambda$, no sobre el costo puro: una solucion mas barata
pero infactible es \emph{peor} desde el punto de vista del
algoritmo.

\subsection{Calculo de $\lambda$ por defecto}\label{sec:lambda}

El helper \code{lambda\_penal\_capacidad\_por\_defecto(ctx)} produce
\begin{equation}\label{eq:lambda}
\lambda^* \;=\; \max\Bigl(\,10\cdot \widetilde{D}_{ij}^{\text{finito}},\;
10\Bigr),
\end{equation}
donde $\widetilde{D}_{ij}^{\text{finito}}$ es la mediana de la matriz
de distancias finitas y positivas de la instancia (precomputada con
Dijkstra). La intuicion es: \emph{una violacion unitaria de capacidad
penaliza tanto como diez deadheads medianos}, lo que en la practica
mantiene la presion sobre la factibilidad sin invalidar movimientos
de reparacion que pasan transitoriamente por infactibles.

% ===================================================================
\section{Operadores de vecindario}\label{sec:operadores}
% ===================================================================

\textsc{MetaCARP} implementa nueve operadores (tres intra y seis
inter) definidos en \code{metacarp/vecindarios.py}. Todos comparten
una firma comun: reciben la solucion en formato normalizado (sin
deposit markers), aplican la transformacion sobre una copia y
devuelven la copia modificada. La eleccion de los indices internos
($i,j,k,l$) es uniformemente aleatoria sobre los valores admisibles.

\subsection{Operadores intra-ruta}

\subsubsection{\op{relocate\_intra}}

\paragraph{Definicion.} Saca la tarea en posicion $i$ de la ruta $r$
y la reinserta en la posicion $j$ ($i\neq j$) de la misma ruta:
\[
R_r = (t_0,\dots,t_{i-1},t_i,t_{i+1},\dots,t_{n-1})
\;\longrightarrow\;
R_r' = (t_0,\dots,t_{i-1},t_{i+1},\dots,t_j,t_i,t_{j+1},\dots,t_{n-1}).
\]

\paragraph{Pseudocodigo.}
\begin{algorithm}[H]
\caption{\op{relocate\_intra}($S, r, i, j$)}
\begin{algorithmic}[1]
\Require Solucion $S$, ruta $r$, posiciones $i\neq j$ con $|R_r|\ge 2$.
\State $S' \gets$ copia profunda de $S$
\State $x \gets S'[r].\text{pop}(i)$
\State $S'[r].\text{insert}(j, x)$
\State \Return $S'$
\end{algorithmic}
\end{algorithm}

\paragraph{Tamano del vecindario.} Para una ruta de longitud $n_r$,
hay $n_r(n_r-1)$ movimientos distintos. Sumando sobre las $m$ rutas
activas se obtiene $\sum_r n_r(n_r-1)$.

\paragraph{Costo computacional.} Un movimiento aleatorio cuesta
$O(n_r)$ por la copia (en la practica, todos los operadores hacen una
copia $O(n)$ donde $n$ es el numero total de tareas). La enumeracion
exhaustiva tiene complejidad $O(n^2)$.

\paragraph{Asignacion tarea $\to$ ruta.} No se altera.

\paragraph{Efecto sobre factibilidad.} \emph{No puede} reducir
violaciones de capacidad: la demanda total por ruta es invariante.

\paragraph{Ejemplo.} $R_1=[A,B,C,D]$ con $i=0$, $j=2$:
$R_1' = [B,C,A,D]$.

\subsubsection{\op{swap\_intra}}

\paragraph{Definicion.} Intercambia las posiciones $i$ y $j$ ($i<j$)
dentro de la misma ruta.

\paragraph{Pseudocodigo.}
\begin{algorithm}[H]
\caption{\op{swap\_intra}($S, r, i, j$)}
\begin{algorithmic}[1]
\Require $|R_r|\ge 2$, $i\neq j$.
\State $S' \gets$ copia profunda de $S$
\State $(S'[r][i], S'[r][j]) \gets (S'[r][j], S'[r][i])$
\State \Return $S'$
\end{algorithmic}
\end{algorithm}

\paragraph{Tamano del vecindario.}
$\sum_r \binom{n_r}{2}$.

\paragraph{Asignacion tarea $\to$ ruta.} No se altera.

\paragraph{Efecto sobre factibilidad.} No reduce violaciones.

\paragraph{Ejemplo.} $R_1=[A,B,C,D]$ con $(i,j)=(0,3)$: $R_1'=[D,B,C,A]$.

\subsubsection{\op{2opt\_intra}}

\paragraph{Definicion.} Revierte el subsegmento $[i,j]$ ($j>i$) dentro
de la ruta $r$.

\paragraph{Pseudocodigo.}
\begin{algorithm}[H]
\caption{\op{2opt\_intra}($S, r, i, j$)}
\begin{algorithmic}[1]
\Require $|R_r|\ge 3$, $0\le i < j < n_r$.
\State $S' \gets$ copia profunda de $S$
\State $S'[r][i\!:\!j+1] \gets \text{reversed}(S'[r][i\!:\!j+1])$
\State \Return $S'$
\end{algorithmic}
\end{algorithm}

\paragraph{Tamano del vecindario.}
$\sum_r \binom{n_r}{2}$.

\paragraph{Efecto.} Heredado del TSP. Particularmente util cuando la
ruta cruza geograficamente: revertir un tramo elimina el cruce.
\emph{No reduce} violaciones de capacidad.

\paragraph{Ejemplo.} $R_1=[A,B,C,D,E]$, $(i,j)=(1,3)$:
$R_1'=[A,D,C,B,E]$.

\subsection{Operadores inter-ruta}

\subsubsection{\op{relocate\_inter}}

\paragraph{Definicion.} Saca la tarea en posicion $i$ de la ruta
$r_a$ y la reinserta en la posicion $j$ de la ruta $r_b\neq r_a$.

\paragraph{Pseudocodigo.}
\begin{algorithm}[H]
\caption{\op{relocate\_inter}($S, r_a, i, r_b, j$)}
\begin{algorithmic}[1]
\Require $|R_{r_a}|\ge 1$, $r_a\neq r_b$.
\State $S' \gets$ copia profunda de $S$
\State $x \gets S'[r_a].\text{pop}(i)$
\State $S'[r_b].\text{insert}(j, x)$
\State \Return $S'$
\end{algorithmic}
\end{algorithm}

\paragraph{Tamano del vecindario.}
$\sum_{a\neq b} n_a (n_b+1)$.

\paragraph{Asignacion.} Transfiere \emph{una} tarea entre rutas.

\paragraph{Efecto sobre factibilidad.} \emph{Si} puede reducir
violacion: si $r_a$ esta sobrecargada y $r_b$ tiene holgura, basta
mover una tarea suficientemente pesada.

\paragraph{Ejemplo.} $R_a=[A,B,C]$, $R_b=[X,Y]$, $(i,j)=(2,1)$:
$R_a'=[A,B]$, $R_b'=[X,C,Y]$.

\subsubsection{\op{swap\_inter}}

\paragraph{Definicion.} Intercambia la tarea en posicion $i$ de
$r_a$ con la tarea en posicion $j$ de $r_b\neq r_a$.

\paragraph{Pseudocodigo.}
\begin{algorithm}[H]
\caption{\op{swap\_inter}($S, r_a, i, r_b, j$)}
\begin{algorithmic}[1]
\Require $|R_{r_a}|\ge 1$, $|R_{r_b}|\ge 1$, $r_a\neq r_b$.
\State $S' \gets$ copia profunda de $S$
\State $(S'[r_a][i], S'[r_b][j]) \gets (S'[r_b][j], S'[r_a][i])$
\State \Return $S'$
\end{algorithmic}
\end{algorithm}

\paragraph{Tamano del vecindario.}
$\sum_{a<b} n_a n_b$.

\paragraph{Efecto sobre factibilidad.} Si las dos tareas tienen
demanda distinta, el balance de capacidad cambia: puede mejorar o
empeorar la violacion. Es un operador \emph{simetrico}: el numero de
tareas por ruta no varia.

\subsubsection{\op{2opt\_star}}

\paragraph{Definicion.} ``Intercambio de colas''. Corta cada ruta en
una posicion $\text{cut}_a$ y $\text{cut}_b$ respectivamente y
permuta las colas:
\[
R_a = H_a \,||\, T_a, \quad R_b = H_b \,||\, T_b
\;\longrightarrow\;
R_a' = H_a \,||\, T_b, \quad R_b' = H_b \,||\, T_a.
\]

\paragraph{Pseudocodigo.}
\begin{algorithm}[H]
\caption{\op{2opt\_star}($S, r_a, c_a, r_b, c_b$)}
\begin{algorithmic}[1]
\Require $r_a\neq r_b$, $0\le c_a < n_{r_a}$, $0\le c_b < n_{r_b}$.
\State $S' \gets$ copia profunda de $S$
\State $H_a, T_a \gets S'[r_a][:c_a{+}1],\; S'[r_a][c_a{+}1:]$
\State $H_b, T_b \gets S'[r_b][:c_b{+}1],\; S'[r_b][c_b{+}1:]$
\State $S'[r_a] \gets H_a + T_b$
\State $S'[r_b] \gets H_b + T_a$
\State \Return $S'$
\end{algorithmic}
\end{algorithm}

\paragraph{Tamano del vecindario.}
$\sum_{a<b} n_a n_b$.

\paragraph{Efecto.} Puede transferir bloques completos de tareas
entre rutas, alterando drasticamente la asignacion. Util para
\emph{deshacer cruces geometricos} entre rutas distintas y para
balancear demanda.

\paragraph{Ejemplo.} $R_a=[P,Q,R,S]$, $R_b=[X,Y,Z]$, $(c_a,c_b)=(1,1)$:
$R_a'=[P,Q,Z]$, $R_b'=[X,Y,R,S]$.

\subsubsection{\op{cross\_exchange}}

\paragraph{Definicion.} Intercambia un segmento contiguo $[i,j]$ de
$r_a$ con un segmento contiguo $[k,l]$ de $r_b\neq r_a$:
\[
R_a' = R_a[:i] \,||\, R_b[k:l{+}1] \,||\, R_a[j{+}1:],\quad
R_b' = R_b[:k] \,||\, R_a[i:j{+}1] \,||\, R_b[l{+}1:].
\]

\paragraph{Pseudocodigo.}
\begin{algorithm}[H]
\caption{\op{cross\_exchange}($S, r_a, i, j, r_b, k, l$)}
\begin{algorithmic}[1]
\Require $|R_{r_a}|\ge 2$, $|R_{r_b}|\ge 2$, $r_a\neq r_b$,
$i<j$, $k<l$.
\State $S' \gets$ copia profunda de $S$
\State $\sigma_a \gets S'[r_a][i:j{+}1]$;\; $\sigma_b \gets S'[r_b][k:l{+}1]$
\State $S'[r_a] \gets S'[r_a][:i] + \sigma_b + S'[r_a][j{+}1:]$
\State $S'[r_b] \gets S'[r_b][:k] + \sigma_a + S'[r_b][l{+}1:]$
\State \Return $S'$
\end{algorithmic}
\end{algorithm}

\paragraph{Tamano del vecindario.}
$O\bigl(\sum_{a<b} n_a^2 n_b^2\bigr)$, dominado por la enumeracion
de los cuatro indices.

\paragraph{Efecto.} Genera vecinos mucho mas disruptivos que un swap
simple: transfiere bloques contiguos manteniendo su orden interno.

\subsubsection{\op{or\_opt\_2} y \op{or\_opt\_3}}

\paragraph{Definicion.} Extrae un bloque contiguo de $k\in\{2,3\}$
tareas de $r_a$ en posicion $i$ y lo reinserta \emph{sin invertir} en
la posicion $j$ de $r_b$. Cuando $r_a=r_b$ (caso degenerado de unica
ruta) se reubica intra evitando $j=i$.

\paragraph{Pseudocodigo (caso $k=2$).}
\begin{algorithm}[H]
\caption{\op{or\_opt\_2}($S, r_a, i, r_b, j$)}
\begin{algorithmic}[1]
\Require $|R_{r_a}|\ge 2$.
\State $S' \gets$ copia profunda de $S$
\State $\beta \gets S'[r_a][i:i{+}2]$
\State eliminar $S'[r_a][i:i{+}2]$
\State insertar $\beta$ en $S'[r_b]$ a partir de la posicion $j$
\State \Return $S'$
\end{algorithmic}
\end{algorithm}

\paragraph{Efecto.} Operadores de \emph{ejection chain} de tamano
fijo. Util para reasignar pares/triples de tareas vecinas que
geograficamente conviene servir juntas.

\subsection{Tabla resumen}

\begin{table}[h]
\centering
\caption{Comparativa estructural de los nueve operadores.}\label{tab:operadores}
\small
\begin{tabular}{lccccc}
\toprule
Operador & Cat. & Tamano vec. (orden) & Costo aleat. & Mueve t.$\to$r. & Reduce viol. \\
\midrule
\op{relocate\_intra}   & intra & $O(n^2)$ & $O(n)$ & no & no \\
\op{swap\_intra}       & intra & $O(n^2)$ & $O(n)$ & no & no \\
\op{2opt\_intra}       & intra & $O(n^2)$ & $O(n)$ & no & no \\
\op{relocate\_inter}   & inter & $O(n^2)$ & $O(n)$ & si & si \\
\op{swap\_inter}       & inter & $O(n^2)$ & $O(n)$ & si & si (parcial) \\
\op{2opt\_star}        & inter & $O(n^2)$ & $O(n)$ & si & si \\
\op{cross\_exchange}   & inter & $O(n^4)$ & $O(n)$ & si & si \\
\op{or\_opt\_2}        & inter & $O(n^2)$ & $O(n)$ & si & si \\
\op{or\_opt\_3}        & inter & $O(n^2)$ & $O(n)$ & si & si \\
\bottomrule
\end{tabular}
\end{table}

% ===================================================================
\section{Metaheuristicas implementadas}\label{sec:metaheuristicas}
% ===================================================================

Las cinco \mh{} se documentan a continuacion siguiendo un esquema
homogeneo: filosofia, estado, criterio de aceptacion, pseudocodigo del
bucle principal, interfaz con el selector de operadores y mecanismos
especiales.

% --- SA -------------------------------------------------------------
\subsection{Recocido Simulado (SA)}\label{sec:sa}

\paragraph{Filosofia.} \mh{} de trayectoria con criterio
probabilistico de aceptacion (Metropolis). Explora intensamente al
inicio (cuando $T$ es alta) y explota localmente al final (cuando $T$
se enfria).

\paragraph{Estado.} Solucion actual $S$, mejor global $S^\star$,
temperatura $T$, contador de niveles sin mejora.

\paragraph{Criterio de aceptacion.} Sea $\Delta=f_\lambda(S')-f_\lambda(S)$:
\begin{equation}
P(\text{aceptar } S') =
\begin{cases}
1, & \Delta\le 0 \\
\exp(-\Delta/T), & \Delta > 0.
\end{cases}
\end{equation}

\paragraph{Parametros instance-aware.}
\begin{itemize}
\item Longitud de cada nivel: $L=n^2$ donde $n=|E_R|$.
\item $T_0 = 20\,d_{\max}/n$;\; $T_{\min}=20\,d_{\max}/n^2$.
\item Factor de enfriamiento $\alpha=0.95$.
\item \emph{Reheat} cuando se acumulan \code{patience}$=50$ niveles
sin mejora, subiendo $T$ a una fraccion \code{reheat\_factor}$=0.5$
de $T_0$.
\item \emph{Kick} inter-ruta cuando se acumulan
\code{max\_iter\_sin\_mejora\_kick} niveles sin mejora (variante
experimental).
\end{itemize}

\paragraph{Pseudocodigo.}
\begin{algorithm}[H]
\caption{Recocido Simulado (con reheat y kick)}
\begin{algorithmic}[1]
\Require Instancia $G$, $\lambda$, $T_0$, $T_{\min}$, $\alpha$, $L$.
\State $S \gets$ mejor solucion inicial; $S^\star \gets S$
\State $T \gets T_0$;\; $\text{nsm}\gets 0$;\; $\text{nsm\_k}\gets 0$
\While{$T \ge T_{\min}$}
  \For{$i = 1 \dots L$}
    \State $(\text{grupo}, v) \gets$ \textsc{SeleccionarGrupo}$(S)$
    \State $S' \gets$ \textsc{GenerarVecino}$(S, \text{grupo})$
    \State $\Delta \gets f_\lambda(S') - f_\lambda(S)$
    \If{$\Delta \le 0$ o $U(0,1) < \exp(-\Delta/T)$}
      \State $S \gets S'$
      \If{$f_\lambda(S) < f_\lambda(S^\star)$} $S^\star \gets S$ \EndIf
    \EndIf
  \EndFor
  \State $T \gets \alpha T$
  \If{nivel sin mejora} $\text{nsm}\!+\!+$; $\text{nsm\_k}\!+\!+$
  \Else{} $\text{nsm}\gets 0$; $\text{nsm\_k}\gets 0$
  \EndIf
  \If{$\text{nsm}\ge \text{patience}$}
     $T \gets T_0\cdot\text{reheat\_factor}$;\;
     reset contadores
  \EndIf
  \If{$\text{nsm\_k}\ge \text{max\_kick}$}
     $S\gets\textsc{KickInter}(S)$;\; reset contadores
  \EndIf
\EndWhile
\State \Return $S^\star$
\end{algorithmic}
\end{algorithm}

\paragraph{Interfaz con el selector.} En cada iteracion del bucle
interno se invoca
\code{seleccionar\_grupo\_operadores\_inter\_intra(rng, viol, ops\_intra,
ops\_inter, ops\_fallback, alpha\_inter, p\_inter)}.
El selector consume \emph{exactamente} un \code{rng.random()} y
devuelve el grupo de operadores a usar (intra o inter) junto con la
bandera de violacion.

% --- TS Simple ------------------------------------------------------
\subsection{Tabu Search Simple (TS Simple)}\label{sec:ts}

\paragraph{Filosofia.} \mh{} de trayectoria \emph{determinista} con
memoria a corto plazo (lista tabu). Acepta \emph{siempre} el mejor
vecino del lote actual aunque empeore, evitando ciclos cortos mediante
la prohibicion temporal de los movimientos recientes.

\paragraph{Estado.} Lista tabu FIFO \code{lista\_tabu} de longitud
fija \code{tabu\_tenure}$=20$, conjunto auxiliar \code{set\_tabu} para
consulta $O(1)$, contador \code{conteo\_tabu} para sincronizacion.

\paragraph{Criterio de aceptacion.} \emph{best non-tabu} con
aspiracion clasica: si todos los vecinos del lote son tabu pero alguno
supera el mejor global conocido, se acepta (la prohibicion se
\emph{aspira}).

\paragraph{Parametros principales.}
\begin{itemize}
\item \code{iteraciones\_max}$=400$, \code{max\_iter\_sin\_mejora}$=100$.
\item \code{tam\_vecindario}$=25$ vecinos por iteracion.
\item \code{tabu\_tenure}$=20$.
\end{itemize}

\paragraph{Pseudocodigo.}
\begin{algorithm}[H]
\caption{Tabu Search Simple}
\begin{algorithmic}[1]
\Require Instancia $G$, $\lambda$, tenure, tam vecindario.
\State $S \gets$ inicial; $S^\star \gets S$; $T_{abu}\gets \emptyset$
\For{$k = 1 \dots K_{\max}$}
  \State $(\text{grupo}, v) \gets$ \textsc{SeleccionarGrupo}$(S)$
  \State $N \gets \{$\textsc{GenerarVecino}$(S, \text{grupo})\}_{i=1}^{|N|}$
  \State Evaluar $f_\lambda$ sobre $N$ por lote vectorizado
  \State $S' \gets \argmin_{S_i \in N \setminus T_{abu}} f_\lambda(S_i)$
  \If{$N \setminus T_{abu} = \emptyset$}
     \State $S' \gets \argmin_{S_i \in N} f_\lambda(S_i)$
  \EndIf
  \If{$f_\lambda(S') < f_\lambda(S^\star)$} \Comment{aspiracion}
     \State $S^\star \gets S'$
  \EndIf
  \State $S \gets S'$;\; $T_{abu}.\text{push}(\text{mov}(S'))$
\EndFor
\State \Return $S^\star$
\end{algorithmic}
\end{algorithm}

\paragraph{Interfaz con el selector.} La eleccion del grupo se hace
\emph{una vez por iteracion} (no por vecino): todo el lote comparte el
mismo modo intra o inter. Esto es coherente con la filosofia de
\emph{best improvement sobre lote homogeneo}.

% --- RTS ------------------------------------------------------------
\subsection{Reactive Tabu Search (RTS)}\label{sec:rts}

\paragraph{Filosofia.} Extension de TS Simple que ajusta
\emph{dinamicamente} la longitud de la lista tabu en funcion de la
deteccion de ciclos.

\paragraph{Estado.} Igual que TS Simple, mas un \code{historial}
hash $\to$ contador de visitas, una variable \code{tabu\_tenure}
que evoluciona en $[\text{min}, \text{max}]$, y contadores
\code{iter\_sin\_repeticion}.

\paragraph{Mecanica reactiva.}
\begin{itemize}
\item Si la solucion actual reaparece en el historial, se \emph{aumenta}
el tenure (multiplicador \code{factor\_aumento}$=1.2$).
\item Si pasan suficientes iteraciones sin repeticion, se \emph{reduce}
(multiplicador \code{factor\_reduccion}$=0.9$).
\item Si una solucion se repite $\ge \text{umbral}=3$ veces, se dispara
un \emph{escape}: limpia lista tabu y aplica
\code{num\_mov\_escape} movimientos aleatorios.
\end{itemize}

\paragraph{Parametros instance-aware.}
\begin{itemize}
\item \code{iteraciones\_max}$\to 20n$; \code{max\_iter\_sin\_mejora}$\to 5n$.
\item \code{tam\_vecindario}$\to \max(20, 2n)$.
\item \code{tabu\_tenure\_inicial}$\to \max(3, \lfloor\sqrt{n}\rfloor)$.
\item \code{tabu\_tenure\_max}$\to \max(10, \lfloor 3\sqrt{n}\rfloor)$.
\end{itemize}

\paragraph{Pseudocodigo (extracto).}
\begin{algorithm}[H]
\caption{Reactive Tabu Search}
\begin{algorithmic}[1]
\State $\tau \gets \tau_{\text{ini}}$;\; $H \gets$ historial vacio
\For{$k = 1 \dots K_{\max}$}
  \State Generar lote, evaluar y elegir $S'$ con aspiracion clasica
  \State Insertar movimiento en \code{lista\_tabu} (capacidad $\tau$)
  \State $h \gets \textsc{Hash}(S')$
  \If{$h \in H$}
     \State $H[h]\!+\!+$;\; $\tau \gets \min(\tau_{\max}, \tau\cdot 1.2)$
     \If{$H[h] \ge 3$}
        \State \textsc{Escape}$(S')$; limpiar $T_{abu}$, limpiar $H$
     \EndIf
  \Else{}
     \State $H[h]\gets 1$
     \If{\text{iter\_sin\_rep} $\ge$ \text{paciencia}}
        \State $\tau\gets\max(\tau_{\min}, \tau\cdot 0.9)$
     \EndIf
  \EndIf
\EndFor
\end{algorithmic}
\end{algorithm}

\paragraph{Interfaz con el selector.} Identica a TS Simple. Durante
el \emph{escape} \emph{no} se aplica el sesgo intra/inter: se usa
seleccion uniforme sobre la lista completa de operadores para no
sesgar la zona de aterrizaje.

% --- ABC ------------------------------------------------------------
\subsection{Artificial Bee Colony Simple (ABC)}\label{sec:abc}

\paragraph{Filosofia.} \mh{} de \emph{poblacion} inspirada en la
busqueda de alimento de las abejas. Tres tipos de individuos:
\begin{itemize}
\item \emph{Empleadas}: explotan localmente cada fuente.
\item \emph{Observadoras}: eligen fuentes por fitness proporcional
(\emph{roulette wheel}).
\item \emph{Scouts}: reemplazan fuentes que llevan demasiados
intentos fallidos.
\end{itemize}

\paragraph{Estado.} $F$ fuentes (soluciones) en formato \emph{IDs},
contadores de abandono por fuente, mejor global.

\paragraph{Parametros instance-aware.}
\begin{itemize}
\item \code{iteraciones}$\to \max(200, 20n)$.
\item \code{num\_fuentes}$\to \max(10, \lfloor 2\sqrt{n}\rfloor)$.
\item \code{limite\_abandono}$\to \max(15, \lfloor n/2\rfloor)$.
\end{itemize}

\paragraph{Pseudocodigo.}
\begin{algorithm}[H]
\caption{Artificial Bee Colony Simple}
\begin{algorithmic}[1]
\State Inicializar $F$ fuentes; evaluar; $S^\star\gets\arg\min$
\For{$k=1\dots K_{\max}$}
  \ForAll{fuente $i$} \Comment{fase empleadas}
     \State $S'_i \gets$ \textsc{GenerarVecino}$(S_i)$
     \If{$f_\lambda(S'_i) < f_\lambda(S_i)$}
        \State $S_i\gets S'_i$; contador$_i\gets 0$
     \Else{} contador$_i$++ \EndIf
  \EndFor
  \ForAll{observadora}
     \State $i\gets$ ruleta $\propto 1/(1+f_\lambda(S_i))$
     \State explotar $S_i$ como en empleadas
  \EndFor
  \ForAll{fuente $i$ con contador$_i\ge\text{limite}$}
     \State $S_i \gets$ \textsc{Aleatoria}; contador$_i\gets 0$
  \EndFor
  \State Actualizar $S^\star$
\EndFor
\State \Return $S^\star$
\end{algorithmic}
\end{algorithm}

\paragraph{Interfaz con el selector.} A diferencia de SA/TS/RTS, ABC
opera en \emph{IDs} y consume el selector via
\code{generar\_vecino\_ids}. El sesgo intra/inter usa la regla
\emph{floor} $P(\text{inter})=\max(p_\text{inter}, 0.8)$ cuando hay
violacion, eliminando \code{alpha\_inter} de la firma.

% --- Cuckoo ---------------------------------------------------------
\subsection{Cuckoo Search}\label{sec:cuckoo}

\paragraph{Filosofia.} \mh{} de poblacion inspirada en el parasitismo
de cria del cuco. Cada iteracion: (i) cada nido genera un
\emph{cuckoo} mediante un \emph{vuelo de Levy discreto}; (ii) el
cuckoo desafia a un nido aleatorio; (iii) se abandonan los peores
nidos con probabilidad $p_a$.

\paragraph{Estado.} $N$ nidos (soluciones) en formato etiquetas,
mejor global.

\paragraph{Vuelo de Levy discreto.} Numero de perturbaciones
encadenadas:
\begin{equation}
n_{\text{pasos}} \;=\; \min\Bigl(12,\; 1 + \bigl\lfloor |N(0,1)|^{1/\beta}\cdot
\text{pasos\_base}\bigr\rfloor\Bigr).
\end{equation}
Cada paso aplica \emph{un} operador del grupo elegido por el selector.

\paragraph{Parametros instance-aware.}
\begin{itemize}
\item \code{iteraciones}$\to \max(200, 20n)$.
\item \code{num\_nidos}$\to \max(10, \lfloor 2\sqrt{n}\rfloor)$.
\item \code{pasos\_levy\_base}$\to \max(3, \lfloor \sqrt{n}/2\rfloor)$.
\item \code{pa\_abandono}$=0.25$, $\beta=1.5$.
\end{itemize}

\paragraph{Pseudocodigo.}
\begin{algorithm}[H]
\caption{Cuckoo Search}
\begin{algorithmic}[1]
\State Inicializar $N$ nidos; $S^\star\gets\arg\min$
\For{$k=1\dots K_{\max}$}
  \ForAll{nido $i$}
    \State $S'_i \gets$ \textsc{VueloLevy}$(S_i, n_\text{pasos}, \text{grupo})$
    \State $j \gets$ aleatorio en $\{1,\dots,N\}$
    \If{$f_\lambda(S'_i) < f_\lambda(S_j)$} $S_j \gets S'_i$ \EndIf
  \EndFor
  \State Abandonar los $\lfloor p_a N\rfloor$ peores; regenerar
  \State Actualizar $S^\star$
\EndFor
\State \Return $S^\star$
\end{algorithmic}
\end{algorithm}

% ===================================================================
\section{Seleccion de operadores: el problema central}\label{sec:seleccion}
% ===================================================================

A continuacion se formalizan las cuatro estrategias de seleccion
ensayadas. Todas comparten la misma interfaz: dado el estado actual
(particularmente la violacion de capacidad) y la lista de operadores
activos, devolver el grupo a usar (o un operador concreto).

\subsection{Seleccion uniforme (baseline)}

Es el comportamiento por defecto de \code{generar\_vecino} cuando se
pasa \code{pesos\_operadores=None} y no se filtra la lista. El
operador se muestrea uniformemente:
\begin{equation}
P(\text{op}_i) \;=\; \frac{1}{|\mathcal{O}|}.
\end{equation}
Ventaja: sin sesgo \emph{a priori}. Desventaja: no aprovecha que las
violaciones de capacidad \emph{solo} se reparan con operadores inter,
por lo que muchas iteraciones se gastan en intras inutiles.

\subsection{Seleccion sesgada \texttt{pesos\_inter\_bias}}

Implementada en \code{pesos\_inter\_bias} y consumida por
\code{seleccionar\_grupo\_operadores\_inter\_intra}. Cuenta
$n_{\text{intra}}=|\mathcal{O}\cap\mathcal{O}_{\text{intra}}|$ y
$n_{\text{inter}}=|\mathcal{O}|-n_{\text{intra}}$. Para cada operador
$o$:

\paragraph{Modo violacion ($\text{viol}>10^{-12}$).}
\begin{equation}
w(o) \;=\;
\begin{cases}
(1-\alpha_{\text{inter}})/n_{\text{intra}}, & o\in\mathcal{O}_{\text{intra}}\\
\alpha_{\text{inter}}/n_{\text{inter}}, & o\in\mathcal{O}_{\text{inter}}.
\end{cases}
\end{equation}

\paragraph{Modo factible ($\text{viol}\le 10^{-12}$).}
\begin{equation}
w(o) \;=\;
\begin{cases}
(1-p_{\text{inter}})/n_{\text{intra}}, & o\in\mathcal{O}_{\text{intra}}\\
p_{\text{inter}}/n_{\text{inter}}, & o\in\mathcal{O}_{\text{inter}}.
\end{cases}
\end{equation}

\paragraph{Variante \emph{Bernoulli + grupo} (la que efectivamente usan SA/TS/RTS).}
Las cinco \mh{} consumen una version equivalente que en lugar de
construir un vector de pesos hace \emph{un} sorteo de Bernoulli para
elegir el \emph{grupo} y luego deja que \code{generar\_vecino} elija
uniformemente dentro del grupo:
\begin{algorithm}[H]
\caption{\textsc{SeleccionarGrupo}$(rng,\text{viol},\mathcal{O}_{\text{intra}},
\mathcal{O}_{\text{inter}})$}
\begin{algorithmic}[1]
\State $p_{\text{ef}} \gets \alpha_{\text{inter}}$ si $\text{viol}>10^{-12}$, si no $p_{\text{inter}}$
\State $u \gets rng.\text{random}()$
\If{$u < p_{\text{ef}}$ y $\mathcal{O}_{\text{inter}}\neq\emptyset$}
   \State \Return $(\mathcal{O}_{\text{inter}}, \text{viol}>10^{-12})$
\EndIf
\State \Return $(\mathcal{O}_{\text{intra}}, \text{viol}>10^{-12})$
\end{algorithmic}
\end{algorithm}
Defaults: $\alpha_{\text{inter}}=0.8$, $p_{\text{inter}}=0.6$.

\subsection{Seleccion binaria estricta}

Implementada en \code{strict\_intra\_inter\_20260524.py}
(\code{seleccionar\_grupo\_strict}). Es \emph{determinista}:
\begin{equation}
\text{grupo} =
\begin{cases}
\mathcal{O}_{\text{inter}}, & \text{si }\text{viol}>10^{-12}\\
\mathcal{O}_{\text{intra}}, & \text{si }\text{viol}\le 10^{-12}.
\end{cases}
\end{equation}
Se conserva el flag de violacion para los contadores, pero la
seleccion \emph{no consume} ningun valor del RNG: es completamente
predecible una vez fijada $\text{viol}$.

Esta variante usa ademas un conjunto \emph{reducido} de cinco
operadores: $\mathcal{O}_{\text{intra}}=\{\op{relocate\_intra},
\op{swap\_intra}\}$ y $\mathcal{O}_{\text{inter}}=\{\op{swap\_inter},
\op{2opt\_star}, \op{cross\_exchange}\}$. Se excluye \op{2opt\_intra}
por considerarlo redundante con \op{swap\_intra} en instancias
pequenas, y \op{relocate\_inter} para favorecer reorganizaciones
simetricas. Esta variante incorpora ademas un \emph{kick reactivo}
inter-ruta cuando se acumulan demasiadas iteraciones sin mejora,
proporcional a $\max(1, \lfloor n/20\rfloor)$ pasos.

\subsection{AOS por Probability Matching}

Implementada en \code{aos\_pm\_20260527.py}. Mantiene un vector de
pesos adaptativo $\mathbf{w}\in[w_{\min},1]^{|\mathcal{O}|}$. La
seleccion del \emph{grupo} sigue siendo binaria estricta (igual que
\code{seleccionar\_grupo\_strict}), pero \emph{dentro del grupo} se
muestrea proporcional a $\mathbf{w}$.

\paragraph{Regla de actualizacion (cada vez que se registra una
mejora global).}
\begin{equation}\label{eq:pm}
w_o \;\gets\; \max\!\Bigl(w_{\min},\; (1-\alpha)\,w_o
+ \alpha\cdot \mathbb{1}[o = o^\star]\Bigr),
\end{equation}
con $w_{\min}=0.05$, $\alpha=0.2$ por defecto. El operador
ganador $o^\star$ recibe \emph{reward} 1 y los demas 0.

\paragraph{Mecanismo de inyeccion.} \code{aplicar\_patch\_aos} hace
tres \emph{monkey-patches}: (i) reemplaza el selector de grupo en el
modulo de la \mh{}, (ii) envuelve \code{generar\_vecino} para inyectar
$\mathbf{w}$ como \code{pesos\_operadores}, y (iii) envuelve
\code{ContadorOperadores.registrar\_mejora} para llamar a
\code{actualizar} cada vez que la \mh{} registra una mejora del mejor
global. El estado AOS se \emph{reinicia} en cada corrida para
trayectorias independientes.

% ===================================================================
\section{Experimentos realizados}\label{sec:experimentos}
% ===================================================================

\subsection{Baseline original (seleccion uniforme)}

\paragraph{Motivacion.} Disponer de una referencia neutra: cada
operador tiene igual oportunidad.

\paragraph{Cambio algoritmico.}
$P(\text{op})=1/|\mathcal{O}|$ uniformemente.

\paragraph{Veredicto.} Aceptable en SA y ABC; en TS Simple y RTS
la baja proporcion de movimientos inter aceptados llevaba a soluciones
factibles pero con \emph{gap} alto. Sirve como punto de partida pero
no es competitivo en instancias con capacidad ajustada.

\subsection{Experimento \texttt{p\_inter\_exp\_2026050524}}

\paragraph{Motivacion.} Probar un sesgo agresivo hacia inter
incluso en estado factible y \emph{anclar} la primera tarea de cada
ruta (para no perder semilla de proximidad al deposito).

\paragraph{Cambio algoritmico.} Variante posicional del generador de
vecinos (\code{vecindarios\_p\_inter\_exp\_2026050524.py}) con
$p_{\text{inter}}=0.8$ permanente y restriccion sobre la posicion $0$
de cada ruta.

\paragraph{Veredicto.} Empeoro 4 de 5 \mh{}. El anclaje rigidizaba el
espacio y el sesgo permanente generaba aceptaciones erraticas en
soluciones que ya eran factibles. Se descarto la idea de \emph{forzar}
posiciones.

\subsection{Experimento \texttt{strict\_intra\_inter\_20260524}}

\paragraph{Motivacion.} Atacar el problema de raiz: si hay violacion,
\emph{nunca} usar intras; si es factible, \emph{nunca} usar inters.
Anadir un kick reactivo para escapar de optimos locales.

\paragraph{Cambio algoritmico.} Selector binario estricto + reduccion
a 5 operadores + kick inter de $\max(1, \lfloor n/20\rfloor)$ pasos
tras cada estancamiento.

\paragraph{Resultados.}
\begin{itemize}
\item ABC: \emph{gap} 41.83\%; factibilidad 100\%.
\item Cuckoo: 40.55\%; 97.1\%.
\item SA: 41.46\%; 100\%.
\item RTS: 30.07\% pero \emph{factibilidad 0\%}.
\item TS Simple: 30.36\% con factibilidad 10.1\%.
\end{itemize}

\paragraph{Veredicto.} El selector estricto mejoro factibilidad en SA,
ABC y Cuckoo. Pero en TS Simple y RTS produjo soluciones de bajo
\emph{gap} pero \emph{infactibles}: la causa raiz fue que ambos
optimizaban costo puro sin penalizar capacidad. Esto motivo el
siguiente experimento.

\subsection{Experimento \texttt{lambda\_grid\_20260525}}

\paragraph{Motivacion.} Si la causa de la infactibilidad en TS/RTS es
$\lambda$ inadecuado, encontrar el $\lambda$ optimo por \mh.

\paragraph{Cambio algoritmico.} \emph{Grid search} sobre
$\lambda_{\text{factor}}\in\{0.5, 1.0, 2.0, 5.0, 10.0\}$, donde
$\lambda=\lambda_{\text{factor}}\cdot\lambda^*$ y $\lambda^*$ es el
default instance-aware de \eqref{eq:lambda}. Se anadio el parametro
\code{lambda\_capacidad} a TS Simple y RTS para que tambien optimizaran
el objetivo penalizado y no el costo puro.

\paragraph{$\lambda$ optimo por \mh.}
\begin{itemize}
\item ABC: 2.0; Cuckoo: 2.0; SA: 0.5; RTS: 0.5; TS Simple: 2.0.
\end{itemize}

\paragraph{Veredicto.} TS y RTS recuperaron 100\% factibilidad; los
\emph{gaps} estabilizaron alrededor de 39--41\%. Este experimento
estableci\'o un nuevo \emph{baseline} de comparacion.

\subsection{Experimento \texttt{aos\_pm\_20260527}}

\paragraph{Motivacion.} Reemplazar la seleccion uniforme dentro del
grupo por seleccion adaptativa segun rendimiento historico.

\paragraph{Cambio algoritmico.} Probability Matching con $\alpha=0.2$,
$w_{\min}=0.05$. \emph{Reward} 1 al operador que registra mejora
global; 0 al resto.

\paragraph{Veredicto.} \emph{No} mejoro significativamente. Tres
hipotesis principales:
\begin{enumerate}
\item Las mejoras globales son \emph{escasas} (decenas por corrida):
con tan pocas actualizaciones, $\mathbf{w}$ apenas se aparta del
uniforme inicial.
\item $\alpha=0.2$ es probablemente demasiado bajo para corridas
cortas (PM es notoriamente sensible a $\alpha$).
\item No hay \emph{warm-start} entre corridas: cada repeticion empieza
con pesos uniformes, perdiendo la memoria acumulada en repeticiones
anteriores.
\end{enumerate}

\subsection{Experimento \texttt{path\_relinking\_20260528}}

\paragraph{Motivacion.} Aprovechar el mejor global como \emph{guia}:
al sufrir un estancamiento, no perturbar al azar sino \emph{redirigir}
la solucion actual hacia el mejor conocido.

\paragraph{Cambio algoritmico.} \emph{Path Relinking} truncado con
representacion \emph{arc-set}:
\begin{enumerate}
\item Calcular el delta = tareas con asignacion (ruta, posicion)
distinta entre la solucion inicio y la guia.
\item En cada paso, mover \emph{una} tarea hacia su posicion guia
eligiendo la transferencia que minimiza $f_\lambda$.
\item Truncar tras $\min(|\text{delta}|, 50)$ pasos.
\item Devolver el \emph{mejor intermedio} (no el final).
\end{enumerate}
Con probabilidad $p_{\text{pr}}=0.5$ se lanza PR despues del kick puro;
con probabilidad $1-p_{\text{pr}}$ se conserva solo el kick. Toda la
captura de \code{mejor\_sol} se hace por \emph{monkey-patch} sobre
\code{copiar\_solucion\_labels} y
\code{ContadorOperadores.registrar\_mejora}, sin tocar las \mh.

\paragraph{Veredicto.} \emph{Primer experimento con mejora consistente}
sobre el mejor $\lambda$ por \mh. Ver Seccion \ref{sec:comparativa}.

% ===================================================================
\section{Comparativa cuantitativa final}\label{sec:comparativa}
% ===================================================================

\subsection{\emph{Gap} medio al \emph{BKS}}

\begin{table}[h]
\centering
\caption{\emph{Gap} medio al \emph{BKS} (\%) por \mh{} y experimento.}
\label{tab:gap}
\small
\begin{tabular}{lcccc}
\toprule
\mh{} & Strict & AOS-PM & $\lambda$-grid best & PR \\
\midrule
SA           & 41.46 & 41.13 & 41.32 & \textbf{39.82} \\
TS Simple    & 30.36 & 40.74 & 39.83 & 39.98 \\
RTS          & 30.07 & 40.21 & 39.44 & \textbf{39.04} \\
ABC          & 41.83 & 41.87 & 41.18 & \textbf{35.80} \\
Cuckoo       & 40.55 & 40.48 & 40.14 & \textbf{36.95} \\
\bottomrule
\end{tabular}
\end{table}

\noindent En negrita los mejores valores por fila. Notese que el
\emph{gap} aparentemente excelente de TS Simple y RTS bajo
\meta{Strict} (30.36\% y 30.07\%) esta inflado: esos valores
corresponden a soluciones mayoritariamente \emph{infactibles}
(ver Tabla \ref{tab:fact}). Tras corregir $\lambda$ los \emph{gaps}
suben a $\sim 39$--$40\%$ pero las soluciones son factibles. La
columna PR muestra mejoras sustanciales en ABC ($-5.38$ pp) y
Cuckoo ($-3.19$ pp).

\subsection{Factibilidad}

\begin{table}[h]
\centering
\caption{Factibilidad final (fraccion de corridas con
$\text{viol}=0$).}\label{tab:fact}
\small
\begin{tabular}{lccc}
\toprule
\mh{} & Strict & AOS-PM & PR \\
\midrule
SA           & 100.0\% & 100.0\% & 100.0\% \\
TS Simple    & 10.1\%  & 100.0\% & 100.0\% \\
RTS          & 0.0\%   & 100.0\% & 100.0\% \\
ABC          & 100.0\% & 100.0\% & 100.0\% \\
Cuckoo       & 97.1\%  & 97.4\%  & 99.1\%  \\
\bottomrule
\end{tabular}
\end{table}

\paragraph{Por que TS y RTS bajo \meta{Strict} casi no eran factibles.}
La causa era estructural: ambos optimizaban \emph{costo puro}
(seleccionaban el vecino con menor costo, no con menor objetivo
penalizado) y la lista tabu no incluia ninguna penalizacion explicita
de la violacion. El selector estricto forzaba el grupo \emph{inter}
solo cuando $\text{viol}>0$, pero como TS aceptaba el ``mejor vecino
del lote'' por costo puro, regresaba inmediatamente al regimen
infactible si una solucion infactible \emph{barata} aparecia en el
lote. El experimento \texttt{lambda\_grid\_20260525} resolvio
esto integrando $\lambda$ en el criterio de seleccion del lote (TS y
RTS pasaron a usar $\text{costo} + \lambda\cdot\text{viol}$).

% ===================================================================
\section{Insights cualitativos}\label{sec:insights}
% ===================================================================

\subsection{Intensificacion vs diversificacion}

Cada estrategia distribuye de forma distinta el esfuerzo entre
explotacion local (intensificacion) y exploracion global
(diversificacion):
\begin{itemize}
\item \meta{Uniforme}: 1/3 intra, 2/3 inter (proporcional al tamano
de los grupos). No reacciona al estado.
\item \meta{Sesgo fluido} (alpha 0.8 / p 0.6): 80\% inter bajo
violacion, 60\% inter en factible. Reacciona al estado.
\item \meta{Strict}: 100\% inter o 100\% intra. Reaccion maxima al
estado pero pierde mezcla.
\item \meta{AOS-PM}: el grupo se elige como en \meta{Strict}; dentro
del grupo, peso proporcional a exito historico. Reaccion al estado +
reaccion al historial.
\item \meta{PR}: ortogonal: no cambia el selector, anade un mecanismo
\emph{guiado} (hacia el mejor global) que se dispara en estancamiento.
Es puramente \emph{intensificador} dirigido.
\end{itemize}

\subsection{Por que PR funciona mejor en \mh{} de poblacion}

ABC y Cuckoo mantienen multiples soluciones en paralelo; el mejor
global cambia frecuentemente (cada vez que un cuckoo o una abeja
genera una mejora). Esto hace que la \emph{guia} de PR sea fresca y
relevante, conduciendo a transferencias informadas. En cambio, en
trayectoria (SA, TS, RTS) el mejor global es estable durante muchos
pasos y PR converge a la guia sin descubrir mejoras intermedias.

\subsection{Por que AOS-PM no mejoro}

Tres factores convergen:
\begin{enumerate}
\item \textbf{Recompensas escasas}: \code{registrar\_mejora} se
dispara solo cuando \emph{se mejora} el mejor global; en
$\sim 100$--$400$ iteraciones esto ocurre en $\le 20$ ocasiones. Con
$\alpha=0.2$, los pesos se mueven muy poco.
\item \textbf{Confusion intra/inter}: el grupo se sigue eligiendo de
forma binaria, asi que el peso de \op{relocate\_intra} se actualiza
solo en iteraciones \emph{intra}; igualmente para los inters. La
serializacion impide la comparacion directa.
\item \textbf{Sin warm-start}: cada corrida arranca con pesos
uniformes, desperdiciando memoria potencial de corridas anteriores.
\end{enumerate}

\subsection{Por que el selector estricto solo es bueno con kick + $\lambda$}

El selector estricto, aislado, es demasiado rigido: si el algoritmo
cae en un optimo local intra, sin movimientos inter no puede salir; si
queda atrapado en una zona infactible, sin movimientos intra no puede
refinar dentro de las rutas. El kick proporciona la salida (perturba
inter), y $\lambda$ adecuado asegura que el criterio de aceptacion
prefiera factibles.

\subsection{El rol del kick reactivo}

Es el unico mecanismo en \textsc{MetaCARP} (aparte del reheat de SA y
el escape de RTS) capaz de \emph{forzar} una salida de un optimo local
sin esperar al criterio termodinamico. Su intensidad,
$\max(1, \lfloor n/20\rfloor)$ pasos, es proporcional al tamano del
problema: pequena en instancias chicas, fuerte en grandes. Sin kick,
el selector estricto y AOS-PM tienden a quedarse en plateaus largos.

% ===================================================================
\section{Conclusiones y trabajo futuro}\label{sec:conclusiones}
% ===================================================================

\subsection{Mejor approach hasta la fecha}

\textbf{Path Relinking truncado con $p_{\text{pr}}=0.5$ sobre el selector
binario estricto y el kick reactivo} es el unico experimento que
mostro mejora consistente en al menos tres de las cinco \mh{} (SA, ABC,
Cuckoo) sin sacrificar factibilidad en ninguna. Para TS Simple y RTS
mostro paridad con $\lambda$-grid (mejora marginal o estabilidad).

\subsection{Recomendaciones de extension}

\begin{itemize}
\item \textbf{Aumentar $p_{\text{pr}}$ en ABC y Cuckoo} (probar 0.7 y
0.9): son las \mh{} donde PR mas brillo, asi que mas PR podria mejorar
mas.
\item \textbf{PR multi-guia}: en lugar de guiar siempre hacia el mejor
global, mantener un pool elite de las 3--5 mejores soluciones distintas
y elegir guia uniformemente entre ellas. Esto diversifica las
direcciones de relinking.
\item \textbf{Pool elite con criterio de diversidad}: agregar
soluciones al pool solo si distan suficientemente (medida arc-set) del
resto, evitando convergencia prematura.
\item \textbf{Combinar PR con AOS}: usar los pesos AOS para elegir,
dentro de PR, \emph{que tarea mover primero} (no solo por menor
$f_\lambda$ sino tambien con preferencia por operadores fuertes).
\item \textbf{Warm-start de AOS-PM entre corridas} con el mismo
seed-de-instancia: salvar el vector $\mathbf{w}$ final en un cache y
recargarlo al inicio.
\end{itemize}

\subsection{Caminos descartados}

\begin{itemize}
\item Forzar la primera posicion de la ruta
(\texttt{p\_inter\_exp\_2026050524}): demasiado rigido,
empeora 4/5 \mh.
\item Optimizar costo puro sin $\lambda$ en TS/RTS: produce
infactibilidad masiva en cualquier variante con selector estricto.
\item Aumentar $\alpha_{\text{aos}}$ por encima de 0.5: hace los
pesos demasiado volatiles y la seleccion degenera a un solo operador.
\end{itemize}

% ===================================================================
\section{Apendice: detalles de implementacion}\label{sec:apendice}
% ===================================================================

\subsection{Monkey-patch de experimentos}

Los experimentos no editan las cinco \mh{} base. Aplican
\emph{monkey-patches} a tres lugares:
\begin{enumerate}
\item \textbf{Selector}: cada modulo de \mh{} (\code{recocido\_simulado},
\code{busqueda\_tabu\_simple}, etc.) importa
\code{seleccionar\_grupo\_operadores\_inter\_intra} en su top-level.
Reasignar el atributo del modulo cambia el selector en uso sin
tocar el archivo.
\item \code{generar\_vecino} y \code{generar\_vecino\_ids}: se envuelven
en una version que inyecta \code{pesos\_operadores} (AOS).
\item \code{ContadorOperadores.registrar\_mejora}: se envuelve para
notificar al AOS / capturar la mejor solucion en PR.
\end{enumerate}

\subsection{Captura de $\lambda$ via \code{sys.\_getframe}}

\code{\_extraer\_ctx\_y\_lambda} sube por la pila de llamadas hasta
encontrar un frame con \code{ctx} entre sus locales; en el mismo
frame busca \code{lam\_eff}, \code{lambda\_eff} o
\code{lambda\_capacidad\_eff}. Si no encuentra ninguno, recalcula
$\lambda^*$ mediante \code{lambda\_penal\_capacidad\_por\_defecto(ctx)}.
Esta tecnica permite a PR conocer la $\lambda$ exacta que la \mh{}
esta usando sin modificar la firma de \code{aplicar\_kick\_*}.

\subsection{Captura de \code{mejor\_sol} via patch de \code{copiar\_solucion\_labels}}

SA, TS Simple, RTS y Cuckoo copian la solucion mediante
\code{copiar\_solucion\_labels(sol)} cada vez que actualizan el mejor
global. PR envuelve esta funcion para que, ademas de hacer la copia,
registre la ultima copia en \code{\_estado\_pr\_local.mejor\_sol\_labels}.
Para ABC (que opera en IDs y no usa este helper), PR envuelve
\code{ContadorOperadores.registrar\_mejora} y extrae
\code{mejor\_sol\_ids} del frame del llamante. El patch es seguro en
la practica porque la unica copia que realmente importa para PR es la
que ocurre dentro del bucle principal de la \mh{}.

% ===================================================================
\begin{thebibliography}{9}

\bibitem{Glover1997}
F.~Glover y M.~Laguna,
\emph{Tabu Search},
Kluwer Academic Publishers, 1997.

\bibitem{Lourenco2003}
H.~R.~Lourenco, O.~C.~Martin y T.~Stutzle,
\emph{Iterated Local Search},
en \emph{Handbook of Metaheuristics}, Springer, 2003, pp. 320-353.

\bibitem{Battiti1994}
R.~Battiti y G.~Tecchiolli,
\emph{The Reactive Tabu Search},
ORSA Journal on Computing 6(2), 1994, pp. 126-140.

\bibitem{Yang2009}
X.-S.~Yang y S.~Deb,
\emph{Cuckoo Search via Levy flights},
en \emph{World Congress on Nature and Biologically Inspired Computing},
IEEE, 2009, pp. 210-214.

\bibitem{Karaboga2005}
D.~Karaboga,
\emph{An idea based on honey bee swarm for numerical optimization},
Technical Report TR06, Erciyes University, 2005.

\bibitem{Fialho2010}
A.~Fialho, L.~da Costa, M.~Schoenauer y M.~Sebag,
\emph{Analyzing bandit-based adaptive operator selection mechanisms},
Annals of Mathematics and Artificial Intelligence 60, 2010,
pp. 25-64.

\bibitem{Resende2005}
M.~G.~C.~Resende y C.~C.~Ribeiro,
\emph{GRASP with Path-Relinking: Recent Advances and Applications},
en \emph{Metaheuristics: Progress as Real Problem Solvers}, Springer,
2005, pp. 29-63.

\end{thebibliography}

\end{document}
